\chapter{Sucesiones y límites}
\label{cap:sucesiones-limites}

Las sucesiones ofrecen una primera formulación precisa de convergencia. Se desarrollan los criterios que permiten detectar límites y se conecta la propiedad de Cauchy con la completitud de $\mathbb{R}$ \cite{abbott2015,rudin1976}.

\section{Sucesiones reales}
Definición, ejemplos, acotación y operaciones término a término.

\section{Definición \texorpdfstring{$\varepsilon$--$N$}{epsilon--N} de límite}
Cuantificadores, prueba directa de límites y unicidad del límite.

\section{Álgebra de límites}
Límites de sumas, productos, cocientes y composiciones cuando corresponda.

\section{Sucesiones monótonas}
Convergencia de sucesiones monótonas y acotadas.

\section{Subsucesiones}
Definición, preservación de límites y uso para probar divergencia.

\section{Teorema de Bolzano--Weierstrass}
Extracción de una subsucesión convergente de toda sucesión acotada.

\section{Límites superior e inferior}
Límites superior e inferior y su uso para describir el comportamiento asintótico.

\section{Sucesiones de Cauchy y completitud}
Criterio de Cauchy, convergencia en $\mathbb{R}$ y equivalencia con la completitud.
